\section{Introduction}

A prophet inequality compares irrevocable online choices with the best
feasible choice after all values are known. Here feasibility is described
by a known finite labeled matroid $M=(E,\mathcal I)$: its independent sets
$\mathcal I$ are closed under taking subsets and satisfy the augmentation
axiom. For a nonnegative vector $v=(v_e:e\in E)$, write
\[
 \OPT_M(v)=\max_{I\in\mathcal I}\sum_{e\in I}v_e.
\]
An online rule observes labels and their values one at a time, and must
accept or reject each label before the next arrival while keeping its
accepted set independent. The value at each label has an unknown
distribution; before arrivals the rule receives one independent sample
from each of these distributions.

We prove a constant guarantee against an \emph{almighty} ordering
adversary. Besides the samples and the entire online value vector, this
adversary sees the rule's complete random seed before choosing the
permutation. The rule is not told that future order. This exposes both
the statistical information and every random choice that might otherwise
protect the online rule from an unfavorable ordering.

\begin{theorem}\label{thm:main}
There is a measurable online selection rule, independent of the value
distributions, with the following property. Let $M=(E,\mathcal I)$ be any
finite known labeled matroid, and let $(S_e,V_e:e\in E)$ be mutually
independent nonnegative real random variables, with $S_e$ and $V_e$ having
the same arbitrary law $D_e$. Assume
$\Ex[\OPT_M(V)]<\infty$. The rule receives exactly the sample vector $S$
before arrivals and uses a seed $R$ independent of $S,V$.

For every measurable permutation rule $\pi$ depending on $M$, the laws
$(D_e)$, $S$, $V$, and $R$, the rule accepts an independent set $A$,
irrevocably deciding at each arrival using only the available history, and
\begin{equation}\label{eq:main}
 \Ex\left[\sum_{e\in A}V_e\right]
       \ge 2^{-310}\Ex[\OPT_M(V)].
\end{equation}
\end{theorem}

The constant is independent of the matroid, its size and rank, and the
distributions. The rule is a finite measurable construction; no
running-time or polynomial independence-oracle guarantee is claimed.
Section~\ref{sec:accounting} proves the stronger constant $2^{-293}$,
leaving slack in the stated bound. We use no distributional description, extra samples, or
randomness concealed from the adversary.

The proof also gives a secretary consequence. In that model the weights
are fixed before any randomness, and the rule observes and rejects a
random initial segment of a uniformly random permutation. After this
prefix, an adversary that sees all weights, the ordered prefix, and the
complete rule seed may rearrange the remaining labels arbitrarily.
Corollary~\ref{cor:secretary} gives a constant guarantee even with this
suffix order. Both results follow from the same fixed-weight guarantee:
after a random set of weights is revealed and its labels sacrificed, the
expected minimum reward over all orders is a constant fraction of the
fixed optimum.

\paragraph{Historical context.}
In the classical single-choice
problem, independent nonnegative values with known distributions admit
a sharp reward fraction of $1/2$; the original work of Krengel and
Sucheston credits Garling for the factor-two bound~\cite{KS77}.
Kleinberg and Weinberg proved the same $1/2$ guarantee for arbitrary matroid
constraints when the distributions are known~\cite{KW12}. Their arrival
model permits the next label to depend on previously revealed values,
without giving the adversary the unseen values.

For unknown distributions, independent labeled samples provide a natural
alternative source of information. Azar, Kleinberg, and Weinberg developed
this limited-information framework and a reduction from order-oblivious
secretary algorithms to single-sample prophet algorithms~\cite{AKW14}.
In the single-choice setting, Rubinstein, Wang, and Weinberg obtained the
optimal one-sample reward fraction $1/2$~\cite{RWW20}. Caramanis et al.
developed greedy-ordered selection and a pointwise sample-pair framework
for single-sample guarantees under structured feasibility
constraints~\cite{CDF22}. These works illustrate how symmetry between a
sample and a realized value can replace knowledge of a distribution.

This connection also links sample-based selection to the matroid
secretary problem, introduced by Babaioff, Immorlica, and
Kleinberg~\cite{BIK07}. In that problem the weights are fixed and the
arrival order is uniformly random. Lachish~\cite{Lachish14} and Feldman,
Svensson, and Zenklusen~\cite{FSZ15} obtained guarantees with a doubly
logarithmic loss in the rank. The latter gave an order-oblivious rule,
whose guarantee survives adversarial reordering after its observation
prefix, and derived a one-sample prophet inequality with the same
loss~\cite[Theorem 1.1 and Corollary 1.2]{FSZ15}. Thus one sample already
gave guarantees for general matroids; we obtain a constant independent
of rank under full-seed order selection.
Recent work also obtains constant guarantees for ordinary random-order
matroid secretary selection~\cite{Singla26,Huang26}.
Section~\ref{sec:secretary} compares these results, which retain uniformly
random arrivals throughout, with the full-seed adversarial-suffix
consequence proved here.

For every fixed $0<\delta<1/4$, Fu et al. obtained a reward fraction
$1/4-\delta$ for arbitrary matroids on $n$ elements, using
$O_\delta(\log^4(2+n))$ samples per element~\cite{FLTWWZ24}.
Feldman, Svensson, and Zenklusen subsequently improved the sample dependence
on the ground-set size and rank~\cite{FSZ26}. These latter results allow
an almighty adversary, which observes all random realizations before
choosing the order, including the algorithm's randomness.
The full-seed adversary is the almighty adversary of the greedy
online-contention-resolution framework of Feldman, Svensson, and
Zenklusen~\cite[Section 1.1]{FSZ16}. The distinction between one sample
and a sample budget growing with the matroid remains essential even when
the latter gives a much larger numerical constant.

For a different arrival model, Abdi, Banihashem, Hajiaghayi, and Mittal
give a one-sample $1/2$ matroid prophet guarantee using $O(n^2)$
independence queries~\cite[Theorem 1.3]{ABHM26}. Their order is fixed
independently of the samples, online values, and internal random coins.
Theorem~\ref{thm:main} instead permits the order to depend jointly on all
of them. A bound for every fixed order controls the minimum of expected
rewards, whereas this information pattern requires an expected minimum
of rewards. The two guarantees therefore have different adversary
interfaces, as well as different quantitative and computational bounds.

The density decompositions used below have a classical foundation in
Edmonds's matroid partition theorem~\cite{Edmonds65}. Soto used principal
partitions and their uniformly dense minors for secretary selection in
the random-assignment model~\cite{Soto13}; Santiago, Sergeev, and Zenklusen
subsequently developed that approach without advance knowledge of the
matroid~\cite{SSZ25}. Our density optimizer is applied within each sampled
weight group and incorporated into paths expanded by independent guards.
The layer construction also has a predecessor in the alternating
restriction--contraction minors along sampled spans used by Feldman,
Svensson, and Zenklusen~\cite[Section 3]{FSZ15}. Here the sampled
density expansions and guards produce the flats, and the analysis must
control the expected minimum over orders after all seed information is
exposed.

\paragraph{Proof strategy.}
We work first with a fixed hidden vector of weights. Independent masks
reveal some coordinates, which are then sacrificed. The goal in this
formulation is the expectation of the \emph{minimum over all orders},
not a guarantee for an order chosen independently of the masks. A simple
coordinatewise coupling subsequently transfers this guarantee to the
one-sample model (Section~\ref{sec:setup}).

After rounding weights into geometric levels, the main rule constructs
nested paths of flats, that is, subsets closed under matroid span. The
path index is an auxiliary integer, unrelated to arrival time. Process
weight groups from low to high. For each group, enlarge the incoming
flats to capture revealed labels at high density relative to rank. A
separate sparse mask of revealed labels supplies \emph{guards}, which
enlarge a flat one auxiliary step before those labels enter its nominal
path. Differences between alternate flats give two-step layers. Within
each layer the rule greedily selects labels independent over the earlier
endpoint flat. Nesting makes their union feasible for every arrival
order (Section~\ref{sec:algorithm}).

The proof must show that these layers retain enough valuable rank. Three
losses arise, each requiring a different argument.

First, lower-weight arrivals may block a focal weight group. At a
revealed label's nominal entry step, the \emph{safety test} asks whether
it remains outside the density expansion of the earlier flat together
with all possible lower-weight competitors. Expose the guards backwards in
auxiliary time, so that the nominal flats shrink. A
disjoint-certificate bound shows that a label usually leaves at a step
with substantial conditional exit probability. A second transition,
using the independent mask that governs online eligibility, turns this
exit estimate into an expected safe count
(Section~\ref{sec:safety}). All masks are drawn initially; reverse
exposure analyzes their law without concealing any coins from the
adversary.

Second, the safe labels themselves have been sacrificed. In each layer,
choose an independent set spanning the unsacrificed labels over the
earlier flat and lower-weight competitors. The density inequality bounds
the safe count by the density charge times the total size of these sets,
plus the number of labels left outside their span. A large safe count can
therefore fail to yield rank only if this residual is large. The residual
lies entirely in the sacrificed mask. Such a set can depend on that mask,
so a bound for one fixed set does not suffice. We encode residuals by
the marks at which their labels first become spanned in two sequences
with boundedly many inserted generators. A combinatorial bound controls
all such marked-pivot patterns before the focal group's mask is exposed.
Crucially, its bound depends on the number of inserted generators, not
the number of auxiliary times. A union bound then rules
out large residuals with high probability, giving rank on unsacrificed
labels (Section~\ref{sec:transfer}).

Third, passing from a group's nominal path to the final guarded path
enlarges the layer bases, which can reduce that rank. One common set of generators accounts for all such
expansions. A telescoping conditional-rank identity charges its cost only
once across the entire path. After independent eligibility thinning,
the remaining rank lower-bounds cumulative accepted counts
simultaneously for every permutation. Only then do we average and sum
over weight levels (Section~\ref{sec:accounting}). A separate
single-choice branch covers the case in which the heaviest label is a
substantial part of the optimum.

Section~\ref{sec:secretary} realizes the sacrificed mask as a random
observation prefix and proves the secretary consequence.
